% chapters/03/section-02.tex
\section{Menyiapkan Perkakas}
\label{sec:menyiapkan-perkakas}

Hanya ada dua perkakas yang wajib ada pada bab ini, yaitu penafsir Python dan
sebuah editor kode. Dua perkakas lain, yaitu terminal dan Git, sangat
dianjurkan karena keduanya akan dipakai terus sepanjang semester.

\subsection{Penafsir Python}

Python adalah bahasa pemrograman yang kodenya dijalankan oleh sebuah program
lain, yaitu \textbf{penafsir} (\istilah{interpreter}). Berkas berisi kode
Python bukan program yang dapat dijalankan langsung oleh sistem operasi; ia
harus diserahkan kepada penafsir. Karena itu, langkah pertama selalu
memastikan penafsir tersedia dan dapat dipanggil.

Sebelum memasang apa pun, periksalah lebih dahulu apakah Python sudah ada pada
komputer yang Anda pakai. Pada komputer laboratorium, Python biasanya sudah
terpasang. Buka terminal, lalu jalankan salah satu perintah berikut.

\begin{kodekeluaran}
python --version
\end{kodekeluaran}

Apabila perintah tersebut tidak dikenali, cobalah bentuk lain yang khusus
tersedia pada Windows.

\begin{kodekeluaran}
py --version
\end{kodekeluaran}

Perintah yang berhasil akan menampilkan versi Python, misalnya
\code{Python 3.11.9}. Buku ini memakai Python 3.8 atau versi yang lebih baru;
versi 3.11 dipakai pada contoh dan tangkapan keluaran di sepanjang buku.

Apabila Python belum ada, unduhlah dari laman resmi Python dan pilih versi
terbaru untuk sistem operasi Anda. Satu hal yang sering terlewat pada
pemasangan di Windows: \textbf{centang pilihan ``Add python.exe to PATH''}
pada layar pertama pemasangan. Tanpa centang tersebut, perintah
\code{python} tidak akan dikenali walaupun Python sudah terpasang, dan Anda
harus mengetikkan alamat lengkap penafsir setiap kali menjalankan program.

\begin{catatabox}
Apabila Anda memakai komputer laboratorium bersama, jangan memasang apa pun
tanpa sepengetahuan pengelola laboratorium. Cukup periksa apakah Python dan
Pygame sudah tersedia dengan menjalankan skrip pemeriksaan pada
Bagian~\ref{sec:memasang-pygame}. Apabila belum tersedia, laporkan kepada
dosen atau pengelola laboratorium.
\end{catatabox}

\subsection{Editor Kode}

Editor kode adalah program untuk menulis kode. Anda boleh memakai editor apa
pun, termasuk editor bawaan yang sederhana, asalkan memenuhi tiga syarat
berikut.

\begin{enumerate}
  \item \textbf{Menyoroti sintaksis}, yaitu mewarnai bagian kode menurut
        jenisnya. Warna berbeda untuk kata kunci, teks, dan komentar membuat
        kesalahan seperti tanda kutip yang belum ditutup mudah terlihat.
  \item \textbf{Mengatur indentasi secara otomatis}. Python sangat peka
        terhadap indentasi, sedangkan editor yang baik membantu Anda
        mempertahankannya secara konsisten.
  \item \textbf{Memudahkan penyimpanan berkas dengan akhiran yang benar},
        yaitu \code{.py}, pada folder yang Anda tentukan.
\end{enumerate}

Editor yang menyediakan terminal terpadu dan penampil pesan kesalahan akan
menghemat waktu, karena Anda tidak perlu berpindah jendela untuk menjalankan
program dan membaca kesalahannya. Pilihan yang lazim dipakai mahasiswa antara
lain Visual Studio Code, PyCharm Community, dan editor teks dengan tambahan
penyorot sintaksis. Buku ini tidak mewajibkan salah satunya; perintah yang
dicontohkan dijalankan dari terminal sehingga berlaku pada editor mana pun.

\subsection{Terminal}

Terminal adalah jendela tempat Anda mengetikkan perintah, termasuk perintah
untuk menjalankan program. Banyak mahasiswa baru menghindari terminal dan
menjalankan program melalui tombol pada editor. Cara itu memang berfungsi,
tetapi terminal tetap perlu dikuasai karena tiga alasan. Pertama, terminal
menampilkan pesan kesalahan secara lengkap, sedangkan sebagian editor hanya
menampilkan sebagian. Kedua, terminal tempat perintah pemasangan paket
dijalankan. Ketiga, terminal dipakai untuk perintah Git pada
Bagian~\ref{sec:struktur-git}.

Hanya dua perintah terminal yang perlu Anda kuasai lebih dahulu, yaitu
berpindah folder dan melihat isi folder.

\begin{kodekeluaran}
cd proyek-pygame\code
dir          (Windows)
ls           (Linux dan macOS)
\end{kodekeluaran}

Perintah \code{cd} berarti \istilah{change directory}, yaitu berpindah ke
folder lain. Perintah \code{dir} atau \code{ls} menampilkan daftar berkas pada
folder yang sedang aktif. Kebiasaan yang perlu dibangun sejak awal adalah
selalu memastikan folder kerja Anda benar sebelum menjalankan program.
Kesalahan ``berkas tidak ditemukan'' yang paling sering terjadi bukan
disebabkan berkasnya hilang, melainkan karena terminal sedang berada di folder
yang keliru.

\subsection{Pengendali Versi Git}

Git mencatat riwayat perubahan berkas. Seluruh tugas praktikum pada mata
kuliah ini dikumpulkan melalui repositori Git pribadi, sehingga Git termasuk
perkakas wajib meskipun bukan untuk menulis kode. Pemasangan dan cara
pemakaiannya diuraikan pada Bagian~\ref{sec:struktur-git}; apabila Git belum
tersedia, pemasangannya dapat ditunda sampai bagian tersebut dibaca.

\begin{catatabox}
Daftar lengkap perkakas beserta kegunaannya pada setiap blok materi dapat
Anda lihat pada Tabel~\ref{tab:perlengkapan} di Bab~I. Yang perlu ada pada
minggu pertama hanyalah Python, editor, dan Git; Unity baru diperlukan mulai
Bab~XII.
\end{catatabox}
